\documentclass[a4paper,11pt]{article}
\usepackage[utf8]{inputenc}
\usepackage{geometry}
\geometry{left = 4.0cm, right = 4.0cm, top = 4.0cm, bottom = 3.5cm}
\usepackage[onehalfspacing]{setspace} 
\usepackage{graphicx}
\usepackage{tabularx}
\usepackage{hyperref}
\usepackage{float} 
\usepackage{ragged2e} %%blocksatz \justifing
\usepackage[ngerman]{babel}
\usepackage[babel,german=guillemets]{csquotes}
\usepackage[backend=biber, style=apa]{biblatex} 
\DeclareLanguageMapping{german}{german-apa}
\newcommand{\showfontsize}{{\f@size pt}}
\addbibresource{literatur.bib}
\setlength{\bibitemsep}{1em}

\title{Abschlusslayout}

\begin{document}
\justifying
\pagenumbering{Roman}



\renewcommand{\contentsname}{Inhaltsverzeichnis}

\tableofcontents\newpage

\newpage
\renewcommand{\listfigurename}{}
\section*{Abbildungsverzeichnis}
\addcontentsline{toc}{section}{Abbildungsverzeichnis}
\listoffigures



\newpage
\addcontentsline{toc}{section}{Tabellenverzeichnis}
\listoftables

\newpage
\section*{Versionen} % change 1: Don't have versions as a numbered section - add toc entry?
\begin{table}[H]
\centering
\begin{tabular}{|p{2.0cm} |p{10.0cm}|}
    \hline
    Version     & Änderungen   \\ \hline
    V1          & Formatvorlage erstellt            \\ \hline
    V2          &  Abstand nach Absätzen geändert, Verzeichnisse an Seitenränder angepasst, Schnellformatvorlagen „Tabelle“ und „Fußnote“ erstellt, Erläuterungen zu Seitenzahlen, Abschnittswechsel und Querverweisen eingefügt               \\ \hline
    V3          & Silbentrennung eingeschaltet, APA-Link ersetzt, Erläuterungen zu Verzeichnissen, Absätzen, Gliederung, Seitenumfang          \\ \hline
    V4          & Konvertierung in Latex \\ \hline
\end{tabular}
\caption[Versionen der Formatvorlage]{Versionen der Formatvorlage}
\end{table}

\newpage
\pagenumbering{arabic}

\section{Einleitung}
\subsection{Das Titelblatt}
Im späteren Verlauf erhalten Sie ein Titelblatt spezifisch für ihren Studiengang. Dieses sollten Sie vor die Seminararbeit anfügen.

\subsection{Formatvorlagen}
Dieses Dokument enthält mehrere Schnellformatvorlagen. Verwenden Sie wenn möglich nur die Formatvorlagen um Ihnen unnötige Formatierungsprobleme zu ersparen. Dies ist insbesondere für die Formatierung von Überschriften wichtig, diese werden sonst nicht im Inhaltsverzeichnis aufgeführt. 

\subsection{Verzeichnisse}
Die Seminararbeit sollte ein Inhaltsverzeichnis, ein Tabellen- und Abbildungsverzeichnis beinhalten – selbstverständlich nur, wenn Tabellen bzw. Abbildungen im Rahmen der Seminararbeit verwendet werden. Sofern Sie die Formatierung beibehalten, können Sie mit Microsoft Word automatisch ein Verzeichnis erstellen lassen. Dies ist auch in OpenOffice möglich. Sie sollten ein Abkürzungsverzeichnis anlegen, sofern Sie Abkürzungen verwenden, die nicht Teil des allgemeinen Sprachgebrauchs sind. Verwenden Sie für die Verzeichnisse jeweils eine eigene Seite.
Außerdem gehört in jede wissenschaftliche Arbeit ein Literaturverzeichnis. Dies steht in der Regel am Ende der Arbeit. Die einzige Ausnahme dessen ist die Verwendung eines Anhangs. Bitte achten Sie bei Referenzen im Text\footnote{Hier steht eine Anmerkung} aber auch im Literaturverzeichnis auf eine einheitliche Zitation und Formatierung (vgl 3.5).







\section{Formatierung}
\subsection{Grundsätzliches}
\subsubsection{Schriftart und Schriftsatz}
Dieser Text ist selbst mit dieser Dokumentvorlage geschrieben und kann in formaler Hinsicht als Muster für die Abfassung von Seminararbeiten verwendet werden. Schriftart Times New Roman oder Times in Schriftgröße 11pt. und Zeilenabstand 1,5. Der Text ist in Blocksatz geschrieben. 

\subsubsection{Gliederung}
Dies ist ein Beispiel für eine dritte Gliederungsebene. Sie sollten in Ihrer Seminararbeit nicht mehr als drei Gliederungsebenen verwenden. Unterkapitel sollten niemals alleine stehen, d.h. wenn Sie eine Unterteilung vornehmen, erstellen Sie immer mindestens zwei Unterkapitel. Die Unterkapitel sollten thematisch und vom Umfang her gleichwertig sind.

\subsubsection{Absätze}
Absätze sind dazu gedacht, unterschiedliche Gedankengänge räumlich voneinander zu trennen und die Lesbarkeit Ihrer Arbeit zu erhöhen. Sie sind nicht geeignet, den Seitenumfang Ihrer Arbeit künstlich aufzublähen. Zu viele Absätze stören außerdem den Lesefluss. Von daher sollten Absätze mindestens vier Zeilen umfassen, dürfen aber auch gerne länger sein. 

\subsection{Seite}
Die Seitenränder lauten wie folgt:
\begin{itemize}
    \item Oben: 4cm
    \item Links: 4cm
    \item Rechts: 4cm
    \item Unten: 3,5cm
\end{itemize}


\subsection{Seitenzahlen}
Die Textseiten Ihrer Arbeit sind mit arabischen Ziffern (1,2, ...) zu nummerieren. Alle anderen Seiten (Inhalts-, Abbildungs-, Tabelle-, Literaturverzeichnisse,  Anhang, etc.) sind mit römischen Ziffern (I, II, ...) zu nummerieren. Dafür benötigen Sie in Word die Funktion des Abschnittswechsels, die in diesem Dokument bereits vorhanden ist. Das Deckblatt zählt als erste Seite, die Seitenzahl wird aber nicht angegeben. Sobald Sie zusätzliche Seiten (wie etwa das Deckblatt) eingefügt haben, müssen Sie ggf. die Nummerierung des Literaturverzeichnisses anpassen. 

\subsection{Abbildungen und Tabellen}
Abbildungen und Tabellen sind durchzunummerieren und mit einem Titel zu versehen (siehe Beispiel Abbildung 1). Außerdem sollten Sie im Text auf das Bild bzw. die Tabelle verweisen. Entweder durch eine Aussage wie bspw. „wie in der folgenden Abbildung zu sehen“ oder aber durch einen Querverweis (s. Abbildung 1). Verwenden Sie dazu die Word-Funktion „Verweise“ $->$ „Beschriftung einfügen“, damit Ihre Verzeichnisse ordnungsgemäß erstellt werden können. Bitte prüfen Sie die Querverweise vor der Abgabe noch einmal besonders sorgfältig, da sie durch Verschiebungen innerhalb des Dokuments schnell kaputtgehen. 
Gleiches gilt für Tabellen:
\begin{table}[H]
\centering
\begin{tabular}{|c|c|c|c|}
    \hline
            & Spalte 1  & Spalte 2  & Spalte 3      \\ \hline
    Zeile 1 & Text      &           &               \\ \hline
    Zeile 2 &           &           &               \\ \hline
\end{tabular}
\caption[Beispieltabelle]{Beispieltabelle}
\end{table}
\noindent Platzieren Sie Ihre Abbildungen und Tabellen möglichst nah an der Stelle, wo sie referenziert werden, aber gleichzeitig platzsparend. Wenn Sie nicht mehr auf die Seite mit der Referenz passen, platzieren Sie sie (wie hier) auf die direkt folgende Seite ganz nach oben und führen Sie den Text auf der vorherigen Seite fort, damit keine halb-leeren Seiten entstehen. 
Abbildungen und Tabellen zählen zum Seitenumfang Ihrer Seminararbeit. Sehr große Abbildungen und Tabellen können im Einzelfall in einen eventuellen Anhang ausgelagert werden, dies sollten Sie individuell mit Ihrem Betreuer besprechen. 

\begin{figure}[H]
    \centering
    \includegraphics[scale=0.7]{infomationssysteme.png}
    \caption[Informationssysteme nach Scheer (Scheer, 2008)] {Informationssysteme nach Scheer (Scheer, 2008)}
\end{figure}
\noindent



\subsection{Zitation}
Zur Zitation von Quellen verwenden Sie bitte den APA-Stil\footnote{http://www.apastyle.org/}. Beispiele für Journalartikel (Lasi, Kemper, Fettke, Feld, \& Hoffmann, 2014), Artikel aus Konferenzbänden (Thaler, Ternis, Fettke, \& Loos, 2015) oder Bücher (Scheer, 2008) finden Sie hier. Weitere Beispiele für unterschiedliche Arten der Litera-tur finden Sie unter anderem \href{https://onlinekurslabor.phil.uni-augsburg.de/course/text/3880/3998}{hier}. Sie können zur Formatierung und Verwaltung der Literatur frei verfügbare Literaturverwaltungsprogramme wie Mende-ley \footnote{https://www.mendeley.com/} oder Zotero\footnote{https://www.zotero.org/}verwenden. Dort sind gängige Zitationsstile wie etwa APA bereits vordefiniert. Über ein Plugin können Sie die Zitationen direkt in Ihr Word-Dokument einfügen. 

\subsection{Fußnoten}
Fußnoten sind nicht für Literaturangaben, sondern nur für zusätzliche Informationen oder Anmerkungen zu verwenden und mit dem Stil \glqq Fußnote\grqq{} (Schriftgröße 10) zu formatieren.


\newpage
\nocite{inproceedings}
\nocite{Scheer1995}
\nocite{Lasi2014}


\pagenumbering{Roman}
\setcounter{page}{5} %CHECK - This might not be right in several cases!!
\renewcommand\refname{Literaturverzeichnis}
\addcontentsline{toc}{section}{Literaturverzeichnis}
%\newpage

\renewcommand*{\mkbibnamefamily}[1]{\textsc{#1}}
\printbibliography
\newpage 
\section*{Eidesstattliche Erklärung}
Hiermit erkläre ich, dass ich die Seminararbeit mit dem Titel

\noindent \rule{\textwidth}{1pt}
eigenständig erbracht, keine anderen als die angegebenen Quellen und Hilfsmittel benutzt und die
aus fremden Quellen direkt oder indirekt übernommenen Gedanken (Texte, Textbausteine und/oder
-fragmente) als solche kenntlich gemacht habe. Die Arbeit wurde nicht, auch nicht in Teilen, unter
Verwendung eines textbasierten Dialogsystems (wie ChatGPT oder andere Werkzeuge basierend auf Large Language Models) oder auf andere Weise mit Hilfe einer
künstlichen Intelligenz von mir verfasst. Die Arbeit habe ich in gleicher oder ähnlicher Form oder
auszugsweise noch keiner Prüfungsbehörde zu Prüfungszwecken vorgelegt.
Des Weiteren bestätige ich, dass die schriftliche und die elektronische Version der Arbeit identisch
sind.
Mir ist bekannt, dass Zuwiderhandlungen gegen den Inhalt dieser Erklärung einen
Täuschungs-versuch darstellen, der grundsätzlich das Nichtbestehen der Prüfung zur Folge hat.

\leavevmode\newline

\noindent \rule{5cm}{1pt}, den \rule{5cm}{1pt}
\leavevmode\newline\newline
\leavevmode\newline
\noindent \rule{10cm}{1pt}


\noindent Unterschrift


\end{document}
